% Source PDF pages 67--69; manual pages 2-38--2-40.
\subsection{Numerical Integration}\label{sec:numerical-integration}

For numerical integration, the equations of motion are written in terms of the state
vector
\begin{equation}
  \vec S(t)
  =\begin{bmatrix}\vec r\\[0.2em]\vec V_{\oplus}\end{bmatrix}
  =\begin{bmatrix}
    \vec r\mathbin{\cdot}\hat x_{\oplus}\\
    \vec r\mathbin{\cdot}\hat y_{\oplus}\\
    \vec r\mathbin{\cdot}\hat z_{\oplus}\\
    \vec V_{\oplus}\mathbin{\cdot}\hat x_{\oplus}\\
    \vec V_{\oplus}\mathbin{\cdot}\hat y_{\oplus}\\
    \vec V_{\oplus}\mathbin{\cdot}\hat z_{\oplus}
  \end{bmatrix}
  =\begin{bmatrix}
    x_{\oplus}\\y_{\oplus}\\z_{\oplus}\\
    \dot x_{\oplus}\\\dot y_{\oplus}\\\dot z_{\oplus}
  \end{bmatrix}.
  \label{eq:trajectory-state-vector}
\end{equation}
The state is a function of time and has a corresponding derivative vector. Given a
time and state,
\begin{equation}
  \dot{\vec S}(t,\vec S)
  =\begin{bmatrix}\dot{\vec r}\\[0.2em]\dot{\vec V}_{\oplus}\end{bmatrix}
  =\begin{bmatrix}\vec V_{\oplus}\\[0.2em]\vec a_{\oplus}\end{bmatrix}
  =\begin{bmatrix}
    \vec V_{\oplus}\mathbin{\cdot}\hat x_{\oplus}\\
    \vec V_{\oplus}\mathbin{\cdot}\hat y_{\oplus}\\
    \vec V_{\oplus}\mathbin{\cdot}\hat z_{\oplus}\\
    \vec a_{\oplus}\mathbin{\cdot}\hat x_{\oplus}\\
    \vec a_{\oplus}\mathbin{\cdot}\hat y_{\oplus}\\
    \vec a_{\oplus}\mathbin{\cdot}\hat z_{\oplus}
  \end{bmatrix}
  =\begin{bmatrix}
    \dot x_{\oplus}\\\dot y_{\oplus}\\\dot z_{\oplus}\\
    \ddot x_{\oplus}\\\ddot y_{\oplus}\\\ddot z_{\oplus}
  \end{bmatrix}.
  \label{eq:trajectory-state-derivative-vector}
\end{equation}
With the earth spin vector aligned with $\hat z_{\oplus}$, the acceleration
components are
\begin{equation}
  \ddot x_{\oplus}
  =\frac{\sum\vec F\mathbin{\cdot}\hat x_{\oplus}}{m}
   +\omega_{\oplus}
      \left(2\dot y_{\oplus}+\omega_{\oplus}x_{\oplus}\right),
  \label{eq:ecfc-x-acceleration-state-form}
\end{equation}
\begin{equation}
  \ddot y_{\oplus}
  =\frac{\sum\vec F\mathbin{\cdot}\hat y_{\oplus}}{m}
   +\omega_{\oplus}
      \left(-2\dot x_{\oplus}+\omega_{\oplus}y_{\oplus}\right),
  \label{eq:ecfc-y-acceleration-state-form}
\end{equation}
\begin{equation}
  \ddot z_{\oplus}
  =\frac{\sum\vec F\mathbin{\cdot}\hat z_{\oplus}}{m},
  \label{eq:ecfc-z-acceleration-state-form}
\end{equation}
where
\begin{equation}
  \omega_{\oplus}=\left\lVert{}_{I}\vec\omega_{\oplus}\right\rVert.
  \label{eq:earth-spin-rate-magnitude}
\end{equation}
The state and derivative vectors are augmented with equations for vehicle mass $m$,
path length $r$, ground range $r_s$, and any user-defined integral variables
(\cref{sec:trajectory-block-define}).

\taossubhead{Vehicle Mass}

The initial mass is part of the trajectory initial conditions
(\cref{sec:block-initial}), and the mass rate is supplied as a function of time and
vehicle state through the propulsive model (\cref{sec:block-prop}). The associated
output variables are
\begin{longtable}{@{}>{\ttfamily}p{0.17\textwidth}p{0.19\textwidth}p{0.56\textwidth}@{}}
  fuel & $\Delta m\,g$ & Fuel used; in English units, $g=32.174\,\mathrm{lb_m/slug}$.\\
  mass & $m$ & Vehicle mass.\\
  mdt  & $\dot m$ & Mass rate.\\
  wt   & $m g$ & Vehicle weight.\\
\end{longtable}

\taossubhead{Path Length}

The path length $r$ is the distance traveled through the air. It is the time integral
of $\lVert\vec V_{\oplus}\rVert$ and is commonly used for launch-rail or sled-track
length. The associated output variables are
\begin{longtable}{@{}>{\ttfamily}p{0.17\textwidth}p{0.19\textwidth}p{0.56\textwidth}@{}}
  plength & $r$ & Path length for the complete trajectory.\\
  plmark  & $r$ & Path length from a designated trajectory point.\\
  plseg   & $r$ & Path length for the current segment.\\
\end{longtable}

\taossubhead{Ground Range}

The ground range $r_s$ is the length of the projection of the flight path onto the
earth's surface. Its derivative is the ground speed
$\lVert\vec V_s\rVert$, the speed of the point on the surface directly below the
vehicle. The ground-speed expression is derived in \cref{sec:ground-speed}.

\taossubhead{Runge--Kutta Integration}

The complete state and derivative vectors evaluated by \texttt{compute\_derivs} are
\begin{equation}
  \vec S(t)
  =\begin{bmatrix}
    x_{\oplus}\\y_{\oplus}\\z_{\oplus}\\
    \dot x_{\oplus}\\\dot y_{\oplus}\\\dot z_{\oplus}\\
    m\\r\\r_s
  \end{bmatrix},
  \qquad
  \dot{\vec S}(t,\vec S)
  =\begin{bmatrix}
    \dot x_{\oplus}\\\dot y_{\oplus}\\\dot z_{\oplus}\\
    \ddot x_{\oplus}\\\ddot y_{\oplus}\\\ddot z_{\oplus}\\
    \dot m\\
    \lVert\vec V_{\oplus}\rVert\\
    \lVert\vec V_s\rVert
  \end{bmatrix}.
  \label{eq:augmented-trajectory-state}
\end{equation}
TAOS integrates these equations with a fixed-step, fourth-order Runge--Kutta
method. Numerical accuracy is therefore controlled by the input step size
$\Delta t$:
\begin{equation}
  \vec S(t+\Delta t)
  =\vec S(t)
   +\frac{1}{6}
    \left(
      \dot{\vec S}
      +2\dot{\vec S}'
      +2\dot{\vec S}''
      +\dot{\vec S}'''
    \right)\Delta t,
  \label{eq:runge-kutta-fourth-order-step}
\end{equation}
where
\begin{equation}
  \dot{\vec S}=\dot{\vec S}(t,\vec S),
  \label{eq:runge-kutta-stage-zero}
\end{equation}
\begin{equation}
  \dot{\vec S}'
  =\dot{\vec S}
    \left(
      t+0.5\Delta t,
      \vec S+0.5\Delta t\,\dot{\vec S}(t)
    \right),
  \label{eq:runge-kutta-stage-one}
\end{equation}
\begin{equation}
  \dot{\vec S}''
  =\dot{\vec S}
    \left(
      t+0.5\Delta t,
      \vec S+0.5\Delta t\,\dot{\vec S}'(t)
    \right),
  \label{eq:runge-kutta-stage-two}
\end{equation}
\begin{equation}
  \dot{\vec S}'''
  =\dot{\vec S}
    \left(
      t+\Delta t,
      \vec S+\Delta t\,\dot{\vec S}''(t)
    \right).
  \label{eq:runge-kutta-stage-three}
\end{equation}
Each integration step requires four derivative evaluations and is performed by
\texttt{integration\_step}. At every time step, TAOS also computes output variables
from the state and derivative vectors, including the longitude, latitude, and
flight-path-angle rates derived below.
